% -------------------------------------------------------------------
% 1. Инструменты и вводная
% -------------------------------------------------------------------

\lecture[17 августа 2026]{1}{Инструменты и~теория}

\textbf{Дисклеймер.} В~этой инструкции описаны основные определения
и~собраны ссылки на~материалы для~обучения в~помощь начинающим
дизайнерам. Все~книги и~курсы авторы читали и~проходили
самостоятельно, поэтому информация предвзята, субъективна
и~не~является актуальным отражением положения дел в~индустрии.
На~момент написания авторы находились на~джуниор"=миддл"=позиции,
поэтому, пожалуйста, относитесь ко~всей информации ниже скептически
и~ищите дополнительные источники информации. Поиск и~проверка
информации~"--- ваш базовый навык как профессионала в~любой сфере,
особенно актуальный во~времена галлюционирующего ИИ~:-)

Дизайн~"--- это всё, что нас окружает. Каждый объект, созданный
человеком: кирпич, учебник, одежда, спутник и~схема города,~"--- это
дизайн. За~любым рукотворным объектом стоят решения и~ход мыслей,
который раньше в~России называли художественным проектированием. 

Дизайн~"--- не~интерфейсы и~не~глянцевые каталоги, он так~же стар,
как и~наша цивилизация. Первые каменные орудия человека уже были
объектами дизайна. Древние люди выбирали подходящие камни,
вытачивали конкретные формы и~получали результат. Это~"---
примитивный дизайн"=процесс, который пройдет долгий путь от~ремесла
к~промышленной революции, мануфактурам, школам дизайна,
индустриализации и~цифровой среде. Чтобы погрузиться в~историю
дизайна, советуем одноимённую книгу А.\,Н.~Лаврентьева.

\begin{definition}[продуктовый дизайнер]
	\tdef{Продуктовый дизайнер}~"--- человек, использующий инструменты
	и~методологию дизайна при~работе над~цифровым продуктом и~влияющий
	на~его~развитие. 
\end{definition}


Продуктовый дизайнер отвечает за то, что~происходит в~продукте:
проектирует новые фичи, реагирует на~обратную связь пользователей
и~бизнеса. UX/UI"=дизайнеры проектируют интерфейсы,
а~продуктовые~"--- исследуют, анализируют, работают с~пользователями
и~стейкхолдерами и~только в~последнюю очередь~"--- проектируют
элементы интерфейса. 

Иными словами, ваша задача как продакта~"--- приносить измеримую
пользу бизнесу и~пользователям. Советуем ознакомиться с~%
\href{https://bureau.ru/soviet/20140623/}{Кодексом бюрошника},
особенно~"--- с~определением полезного действия. Каждое ваше
дизайнерское решение приближает его~достижение. Когда вы понимаете,
какую задачу решаете, для кого, зачем и~каков её результат, вам
проще защищать свои решения перед командой и~бизнесом. Вы всегда
должны быть готовы ответить на~вопрос «почему в~продукте было
сделано именно~так». Хороший продуктовый дизайнер знает свой продукт
и~пользуется им, знаком с~продуктами конкурентов и~следит
за~трендами в~индустрии.


Ключевой термин в~работе продуктового дизайнера~"--- \tdef{метрики},
или~количественные показатели продукта. Существует большое
количество метрик для~оценки состояния продукта~"--- например, MAU,
DAU, Retention, CTR. Каждая метрика прямо или~косвенно отвечает
на~вопрос «что происходит с~продуктом».

Не все школы дизайна считают, что~продуктовый дизайнер обязан
погружаться в~метрики. Тем не~менее полезно знать хотя~бы некоторые
из~них, чтобы говорить на~одном языке с~менеджерами продукта
и~стейкхолдерами~"--- владельцами бизнеса. 

\subsection{Системы обучения}

Всё перечисленное выше можно и~нужно уложить в~единую систему
обучения. Рекомендуем воспользоваться готовой системой.


\begin{enumerate}
	\item
		\href{%
			https://bureau.ru/school/designers%
		}{%
			Школа дизайнеров (Бюро Горбунова)%
		}~"--- фундаментальные знания о~типографике, интерфейсах,
		редактуре и~переговорам. Советуем изучить программу школы
		(см.~\picref{pic:bureau_program})
		и~список рекомендуемых книг. Учёба потребует сил, но~подойдёт
		и~новичкам, и~практикующим дизайнерам. В~школе есть бесплатные
		места, на~которые поступают по~конкурсу. Успешно закончившим все
		три ступени могут предложить работу в~Бюро
		и~в~компаниях"=партнёрах.


	\begin{figure}[h]
		\centering
		\includegraphics[width=\textwidth]{program.jpg} 
		\caption{Дисциплины и~темы, которые проходят в~Школе дизайнеров}
		\label{pic:bureau_program}
	\end{figure}

	\item
		\href{%
			https://intuition.team/ru/design-roadmap%
		}{%
			Как стать дизайнером (Женя Арутюнов, Интуиция)%
		}~"--- список тем с~порядком их~изучения. Подойдёт для учёбы
		в~своём темпе и~для~углубления знаний по конкретным вопросам.
		Если не уверены в~интересе к~дизайну, советуем начать отсюда.

	
	% С помощью \- явно задаём LaTeX'у места, где можно переносить
	% слово
	\item \href{%
			https://wannabe.ru/design/superpowered%
		}{%
			Superpowered UX/UI (cтудия PINKMAN)%
		}~"---курс, который постоянно обновляет программу
		и~инс\-т\-рументы. Не подойдёт, если вы только начали работать
		в~Фигме или сомневаетесь в~своём направлении. Плюс~"---
		актуальные знания. Вы~будете пользоваться последними
		инструментами, знать, как работают практикующие дизайнеры,
		и~сможете выстроить дизайн"=процесс. 

	\item
		\href{%
			https://www.duosapiens.ru/%
		}{%
			Менторство (Duo Sapiens)%
		}~"--- регулярные встречи с~опытным дизайнером. Ментор поможет
		найти пробелы в~знаниях, решить рабочий вопрос или~составить
		резюме. Менторы сами нанимают дизайнеров и~знают, с~какими
		задачами вы~столкнётесь. Рекомендуем, если знания есть,
		но~непонятно, куда двигаться дальше.

\end{enumerate}


\subsection{Книги}

% Замени mod на orig, если хочется вернуть оригинальные картинки
\noindent
\bookitem{books/mod/typo.png}{Типографика и~вёрстка}{Артём Горбунов}\hfill
\bookitem{books/mod/ui.png}{Пользовательский интерфейс}{Илья Бирман}\hfill
\bookitem{books/mod/grid.png}{%
	Модульные системы в~графическом дизайне%
}{Йозеф Мюллер-Брокманн}
\par
\vspace{1.5em}
\noindent
\bookitem{books/mod/ui2.png}{%
	Интерфейс: новые направления в~проектировании компьютерных систем%
}{Джеф Раскин}\hfill
\bookitem{books/mod/font.png}{О~шрифте}{Эрик Шпикерманн}\hfill
\bookitem{books/mod/kovodstvo.png}{Ководство}{Артемий Лебедев}



% Выводы
% \begin{takeaways}
%		\begin{itemize}
%			\item Вывод «А»;
%			\item Вывод «Б»;
%			\item Вывод «В»;
%		\end{itemize}
% \end{takeaways}

% vim:ts=2:sw=2:tw=68
